\section{Related work}\label{sec:related-work}

The relationship between regression discontinuity with measurement error and the deconvolution literature depends on the assignment mechanism, the target, and the criterion used to evaluate statistical procedures. Our \leanref{S-1}{observed-treatment experiment} records the side of the latent cutoff and a Gaussian-contaminated score. The benchmark in \cref{obj:def:model} combines known compact support, quantitative density bounds, binary outcomes, and a fixed Hölder bound on the potential-outcome means. Its \leanref{S-2}{decision problem} evaluates estimators by worst-case expected absolute error and connected closed confidence intervals by worst-case expected length, subject to coverage of at least \(0.9\) for every benchmark law at each sample size. Three comparisons clarify how these criteria relate to existing results.

\subsection{Identification and cutoff decision criteria}

The closest identification comparison is \citet[Assumptions 1Y, 2C, and 7; Proposition 4(b)]{PeiShen2017DevilTails}. Their result identifies the sharp regression discontinuity effect from the observed outcome, treatment status, and mismeasured assignment variable under nondifferential measurement error, positive latent-score density near the cutoff, and Gaussian error. Observed treatment supplies information about the latent cutoff side, while Gaussian convolution permits recovery of the relevant latent distributions. Identification concerns recovery from the population observable law. The present benchmark evaluates the precision available from a finite sample under explicit quantitative restrictions. In particular, the known noise scale and the density and smoothness conditions in \cref{obj:def:model} specify the class over which estimation risk and interval length are compared.
% [Pei and Shen, Proposition 4(b)](https://arxiv.org/pdf/1609.01396)

A complementary inference comparison comes from \citet[Chapter~2, Assumptions~8, 9, 11, and 13; Theorem~6; printed pp.~49, 59, and 61]{Dong2018Thesis}. For observed treatment and a known error density, Theorem~6 establishes asymptotic normality of the deconvolution local-linear cutoff estimator after centering by its bias and standardizing by its variance. For Gaussian error, the prescribed bandwidth has order \((\log n)^{-1/2}\). The theorem also imposes kernel and transform-integrability conditions and a Lyapunov condition; Assumption~11 includes twice differentiable latent density and one-sided regression smoothness near the cutoff. These conditions support a pointwise distributional approximation for local-linear inference. Our criteria instead place the supremum over the specified benchmark class inside the risk and coverage requirements. The finite-sample guarantee in \cref{obj:thm:finite-certificate} and the rate comparison in \cref{obj:thm:uniform-frontier} address these uniform criteria under the benchmark's density bounds and Hölder restrictions. This theorem-level comparison relies on the cited-source attestation; its provenance limitation is recorded in the verification appendix, and it has no proof-use in our results.
% [Hao Dong, Essays in microeconometrics, Chapter 2](https://etheses.lse.ac.uk/3756/1/Dong__essays-in-microeconometrics.pdf)

Assignment and target distinctions also matter within the regression discontinuity literature. \citet{Eckles2025} exploit known exogenous noise in the score used to assign treatment. Their approach uses the resulting assignment randomness to estimate weighted treatment effects associated with changes in the assignment rule. In our experiment, assignment follows the latent score, and measurement error enters the score available to the researcher. The target is the potential-outcome contrast at the latent cutoff. Auxiliary information provides another route: \citet{DaveziesLeBarbanchon2017ContinuousError} study a two-sided fuzzy design with nondifferential measurement error and an auxiliary sample of treated individuals containing both true and noisy scores. Their target is a local average treatment effect, and their auxiliary observations identify features of the measurement process.
% [Eckles et al., Noise-induced randomization](https://arxiv.org/html/2004.09458)
% [Davezies and Le Barbanchon, Continuous measurement error](https://docs.iza.org/dp10801.pdf)

Observed-score interpretations supply a further comparison. Under correct classification of assignment and continuity of the potential-outcome means conditional on the observed score, \citet[Section~2.1, Lemma~2.1]{DongKolesar2023IgnoreMeasurementError} identify the treatment effect for units at the observed-score cutoff. Their additional condition on treatment-effect dependence on measurement error gives a weighted-average interpretation over latent scores. These conditions align conventional regression discontinuity analysis with an observed-score target. The latent-cutoff benchmark specifies a different conditioning event and uses the recorded treatment side together with the proxy distribution to recover its endpoint contrast.
% [Dong and Kolesár, Section 2.1](https://www.princeton.edu/~mkolesar/papers/rd_rounded.pdf)

\subsection{Compact support, boundary density estimation, and moments}

Compact-support deconvolution provides the second decision-problem comparison. \citet[Theorem~1(a), fixed-support case]{Meister2007CompactSupport} obtain uniform mean integrated squared error of order \((\log n/\log\log n)^{-2s}\) for density estimation over a compactly supported Fourier--Sobolev class of smoothness \(s\). Their procedure uses a known bound on the support and knowledge of the error characteristic function on a neighborhood of zero, where it is bounded away from zero. Thus, compact support supplies information that can improve density recovery even under partial knowledge of the error distribution. Our benchmark uses known support and a fully specified Gaussian measurement process, with Hölder smoothness imposed on conditional outcome means. The loss concerns an endpoint causal contrast, while the score density enters as a nuisance subject to continuity and quantitative bounds. These rates characterize different decision problems: integrated density recovery under Fourier--Sobolev restrictions and endpoint contrast recovery under conditional-mean Hölder restrictions.
% [Meister, Theorem 1(a)](https://www.f08.uni-stuttgart.de/.content/media/downloads/Mathematik/mathematische_berichte/2005/2005-007.pdf)

Boundary deconvolution and Gaussian moment methods clarify the constructional relationship. \citet{ZhangKarunamuni2009Boundary} study a density supported on a half-line under known independent additive error. They construct boundary deconvolution kernels and analyze pointwise mean squared error for density and density-derivative estimation under ordinary-smooth and supersmooth errors; their construction can also be adapted to compact support. This is a direct endpoint antecedent, but its target is the latent density rather than a marked-to-unmarked regression ratio, and its criterion does not include honest confidence-interval length. \citet[Theorem~1; Section~4.1]{WuYang2020DenoisedMoments} estimate finite Gaussian mixing distributions under Wasserstein loss, using bounded component locations and a specified number of components; their theorem distinguishes known and unknown common variance. For known variance, Hermite polynomials supply unbiased latent-moment estimates, which their procedure projects onto a feasible moment space before reconstructing the mixing distribution. The finite inverse-heat transform in \cref{obj:def:inverse-heat} uses the related polynomial inversion principle to estimate endpoint-weighted averages. The arm ratios in \cref{obj:def:ratio} then convert marked and unmarked averages into conditional-mean estimates. This connects Gaussian moment inversion to the cutoff target through an explicitly defined observable procedure.
% [Zhang and Karunamuni, Boundary deconvolution](https://doi.org/10.1016/j.jspi.2008.10.021)
% [Wu and Yang, Theorem 1 and Section 4.1](https://arxiv.org/html/1807.07237)

Deconvolution regression gives a closer antecedent for the ratio structure. \citet[Section~2, Theorems~2 and 5--6]{Fan1993} observe a response together with an error-contaminated covariate and estimate the latent regression function by normalizing a deconvolved response numerator by a deconvolved density denominator. They assume a known error law with nonvanishing characteristic function, a latent density bounded away from zero on a compact interval, and smoothness of both the density and regression function. For a supersmooth error of order \(\beta_e\), their uniform squared-risk upper rate is \(O((\log n)^{-2k/\beta_e})\); Gaussian error has \(\beta_e=2\), yielding the displayed convergence rate \((\log n)^{-k/2}\) in root mean square, with matching minimax results under their stated classes. That result concerns regression at points in a fixed interval with a fixed error law. The present target combines two recorded treatment arms at the latent support endpoint and evaluates expected absolute error and honest interval length uniformly over the full range of \(\sigma\).
% [Fan and Truong, Sections 2--3](https://archive.ymsc.tsinghua.edu.cn/pacm_download/468/10828-Jianqing_Fan_37.pdf)

Convolution-model linear functionals provide a complementary decision-theoretic comparison. \citet[Theorem~3.1 and Section~4.1]{Butucea2009} estimate \(\langle\psi,g\rangle\) from noisy observations under a known convolution density whose Fourier transform is nowhere zero. Their adaptive oracle bound balances the squared Fourier-tail bias against a cutoff-dependent penalty, and their pointwise-density specialization recovers the usual pointwise deconvolution regimes. \citet[Theorems~1--2 and Section~2.5]{Pensky2017} give upper and lower quadratic-risk bounds for general linear functionals over Sobolev signal classes with ordinary- or supersmooth Fourier behavior; the bounds match up to constants in the cases identified there and otherwise may differ by a logarithmic factor. Pointwise density is one of their applications. These results locate the inverse-problem difficulty for a single linear functional. Our endpoint estimand is instead assembled from four such marked and unmarked quantities through stabilized ratios, and the benchmark asks simultaneously for all-noise-scale absolute-risk and honest-length guarantees.
% [Butucea and Comte, Theorem 3.1 and Section 4.1](https://helios2.mi.parisdescartes.fr/~comte/ButuceaComte.pdf)
% [Pensky, Theorems 1--2 and Section 2.5](https://arxiv.org/pdf/1411.1660)

\subsection{Vanishing noise and the estimation target}

Shrinking-noise density estimation supplies the third comparison. \citet[Theorems~2.1 and 2.2]{Delaigle2008ShrinkingNoise} analyze mean integrated squared error for deconvolution density estimators as sample size increases and measurement-error scale decreases. Under their density-smoothness, kernel, and error assumptions, sufficiently small noise permits the corresponding error-free density rate; larger noise introduces a scale-dependent cost. This makes the comparison between noise scale and statistical resolution central to density estimation. For the present endpoint criterion, \cref{obj:def:uniform-rate,obj:thm:uniform-frontier} give the corresponding rate characterization across the benchmark noise scales, with constants uniform in sample size and noise scale at fixed benchmark smoothness.
% [Delaigle, Theorems 2.1 and 2.2](https://www3.stat.sinica.edu.tw/statistica/oldpdf/A18n311.pdf)

A recent minimax comparison is \citet[Section~2, Theorems~2.4--2.6]{DurotMukherjee2026VanishingNoise}. Their deconvolution analysis estimates a signal distribution under Wasserstein-one loss over a moment-bounded class, with known error distribution and known noise scale. Under Gaussian noise with standard deviation \(\sigma_n\), Theorem~2.4 gives the lower bound order \(\sigma_n(\log n)^{-1/2}+n^{-1/2}\). Their upper bounds match when \(\sigma_n\ge n^{-1/2+\eta}\), for fixed \(\eta>0\), and when \(\sigma_n\le n^{-1/2}\); the intervening regime carries separately stated logarithmic bounds under additional reciprocal-characteristic-function regularity. Our comparison concerns absolute error in a latent-cutoff contrast and expected length under uniform finite-sample coverage, so these distribution-recovery rates provide a regime comparison rather than a numerical benchmark for our loss.
% [Durot and Mukherjee, Section 2](https://arxiv.org/pdf/2404.09306)

\subsection{Honest inference under smoothness restrictions}

The interval criterion belongs to the optimal-recovery and smoothness-based regression literature. \citet{Donoho1994} connect statistical estimation over restricted function classes with optimal recovery. \citet{ArmstrongKolesar2018OptimalInference} derive finite-sample optimal confidence procedures for linear functionals in regression models with normal errors and known variance, together with uniform asymptotic extensions. \citet{ArmstrongKolesar2020SimpleHonest} construct honest intervals by incorporating worst-case smoothing bias into the critical value and bandwidth choice. Our finite-sample certificate similarly combines a public bias bound with a stochastic-error bound, here for inverse-heat averages and stabilized arm ratios. Coverage is imposed through \cref{obj:def:honest-class}, and precision is evaluated by worst expected length over that same class. These choices give the endpoint estimation and inference results a common decision-theoretic criterion.
% [Donoho, Statistical estimation and optimal recovery](https://digicoll.lib.berkeley.edu/record/85850/files/214.pdf)
% [Armstrong and Kolesár, Optimal inference](https://arxiv.org/abs/1511.06028)
% [Armstrong and Kolesár, Simple and honest confidence intervals](https://doi.org/10.3982/QE1199)
